%=========== ΛΥΣΕΙΣ - ΔΙΑΓΩΝΙΣΜΑ ΠΡΟΣΟΜΟΙΩΣΗΣ 12 ===========

%=========== ΛΥΣΗ - 12ο ΘΕΜΑ Α (Σ/Λ μόνο) ===========
\begin{thema}{Α}
\begin{erwthma}
\setcounter{enumi}{3}
\item \ \\
\begin{alist}
\item \textbf{Λ.} 
Το Θεώρημα \eng{Bolzano} εξασφαλίζει την ύπαρξη \textbf{τουλάχιστον μίας} ρίζας. Δεν αποκλείει, για παράδειγμα, η $f$ να έχει επιστρέψει στον άξονα των $x$ πολλές φορές (άρα η εξίσωση $f(x)=0$ να έχει περισσότερες, συγκεκριμένα μονό πλήθος ριζών). Μόνο αν συμπληρωθεί από γνήσια μονοτονία εγγυάται ακριβώς μία.
\item \textbf{Σ.}
Πρόκειται για το Θεμελιώδες Θεώρημα του Ολοκληρωτικού Λογισμού: η ολοκλήρωση της παραγώγου αποδίδει τη μεταβολή της αρχικής συνάρτησης.
\item \textbf{Λ.}
Αν η συνάρτηση είναι κοίλη, η εφαπτομένη της βρίσκεται (με εξαίρεση το σημείο επαφής) \textbf{πάνω} από τη γραφική παράσταση $C_f$, όχι κάτω.
\item \textbf{Σ.}
Αφού είναι πολυώνυμο άρτιου πλήθους συντελεστών με μεγαλύτερο εκθέτη περιττό $3$, τα όριά της στο $+\infty$ και $-\infty$ δίνουν πάντα διαφορετικά πρόσημα ($+\infty$ και $-\infty$ ή αντιστρόφως). Εφαρμόζοντας το θεώρημα \eng{Bolzano} σε αρκετά μεγάλο διάστημα, εξάγεται η βεβαιότητα ύπαρξης μίας τουλάχιστον πραγματικής ρίζας.
\item \textbf{Λ.}
Απαιτείται επιπλέον η $f''(x)$ να \textbf{αλλάζει υποχρεωτικά πρόσημο} εκατέρωθεν του $x_0$ για να θεωρηθεί το σημείο $(x_0, f(x_0))$ σημείο καμπής.
\end{alist}
\end{erwthma}
\end{thema}

%=========== ΛΥΣΗ - 12ο ΘΕΜΑ Β ===========
\begin{thema}{Β}
Η $f$ δίνεται από το γράφημα (ημικύκλιο $\rho=2$ στο $[-2, 2]$ και ευθύγραμμο τμήμα στο $[2, 5]$).
\begin{erwthma}

\item \textbf{Σύνολο Τιμών:}
Στο εύρος $[-2, 2]$ των τετμημένων, το ημικυκλικό τόξο έχει ως ελάχιστο ύψος το $y=0$ και ως μέγιστο ύψος (την ακτίνα) το $y=2$. Άρα προσφέρει το σύνολο τιμών $[0, 2]$.
Στο εύρος $[2, 5]$, το ευθύγραμμο τμήμα ξεκινά από το $(2,0)$ έως το $(5, 3)$. Εφόσον είναι συνεχές, λαμβάνει κάθε τιμή στο εύρος $[0, 3]$.
Η ένωση των δύο συνόλων $[0, 2] \cup [0, 3] = [0, 3]$, οπότε το συνολικό σύνολο τιμών της $f$ είναι το διαστημα $\boldsymbol{[0, 3]}$.

\item \textbf{Υπολογισμός Εμβαδών Γεωμετρικά:}
- Το $\dint_{-2}^2 f(x)\,\d x$ ταυτίζεται με το εμβαδόν του ημικυκλίου ακτίνας $\rho=2$.
Άρα $E_1 = \frac{1}{2} \cdot \pi\rho^2 = \frac{1}{2} \cdot \pi (2^2) = \boldsymbol{2\pi \ \text{τ.μ.}}$

- Το $\dint_2^5 f(x)\,\d x$ ταυτίζεται με το εμβαδόν του ορθογωνίου τριγώνου, το οποίο σχηματίζεται με βάση $(5-2)=3$ και ύψος του σημείου $(5, 3)$ δηλαδή $3$.
Άρα $E_2 = \frac{1}{2} \cdot 3 \cdot 3 = \boldsymbol{\frac{9}{2} \ \text{τ.μ.}}$

\item \textbf{Παραγωγισιμότητα στο $x=2$:}
Ελέγχουμε τα πλευρικά όρια του λόγου μεταβολής.
Στο διάστημα $[2, 5]$ ισχύει η ευθεία που διέρχεται από τα $(2,0)$ και $(5,3)$. Η κλίση της είναι $m = \frac{3-0}{5-2} = 1$. Άρα ο τύπος της είναι $f(x) = 1(x-2) \Rightarrow f(x) = x-2$.
Η δεξιά παράγωγος στο $x=2$ είναι $\lim_{x\to 2^+} \frac{f(x)-f(2)}{x-2} = \lim_{x\to 2^+} \frac{x-2-0}{x-2} = \boldsymbol{1}$.

Από τα αριστερά, η $f$ περιγράφει το ημικύκλιο $\rho=2$: $x^2 + y^2 = 4 \Rightarrow y = \sqrt{4-x^2}$.
Η αριστερή παράγωγος αποδίδεται από την εφαπτομένη στον κύκλο στο εξωτερικό του $(2,0)$, η οποία γραφικά είναι \textbf{κατακόρυφη} (παράλληλη στον άξονα $y'y$). Αναλυτικά:
\[ \lim_{x\to 2^-} f'(x) = \lim_{x\to 2^-} \frac{-x}{\sqrt{4-x^2}} = \frac{-2}{0^+} = \boldsymbol{-\infty} \]
Εφόσον η αριστερή παράγωγος δεν είναι πεπερασμένος αριθμός (και δεν ταυτίζεται με τη δεξιά) η $f$ \textbf{δεν είναι παραγωγίσιμη} στο $x=2$. 

\item \textbf{Εξίσωση Εφαπτομένης στο $(0, 2)$:}
Το σημείο $(0, 2)$ βρίσκεται ακριβώς στην κορυφή του ημικυκλίου. Γεωμετρικά, η εφαπτομένη σε αυτό το σημείο είναι απολύτως \textbf{οριζόντια} (παράλληλη στον άξονα των τετμημένων).
Κατά συνέπεια ο συντελεστής διεύθυνσης λαμβάνει μηδενική κλίση $\lambda = f'(0) = 0$.
Η εξίσωση γράφεται $y - 2 = 0 \cdot (x - 0) \Rightarrow \boldsymbol{y = 2}$.

\end{erwthma}
\end{thema}

%=========== ΛΥΣΗ - 12ο ΘΕΜΑ Γ ===========
\begin{thema}{Γ}
Δίνεται η συνάρτηση $f(x) = \ln(x^2+1) - x + 1, x\in\mathbb{R}$.
\begin{erwthma}

\item \textbf{Μονοτονία και Ακρότατα:}
Παραγωγίζουμε χρησιμοποιώντας τον κανόνα αλυσίδας:
\[ f'(x) = \frac{1}{x^2+1}(x^2+1)' - 1 = \frac{2x}{x^2+1} - 1 \]
Μορφοποιούμε το κλάσμα με ομώνυμα:
\[ f'(x) = \frac{2x - (x^2+1)}{x^2+1} = \frac{-x^2 + 2x - 1}{x^2+1} = -\frac{x^2 - 2x + 1}{x^2+1} = -\frac{(x-1)^2}{x^2+1} \]

Διότι το τετράγωνο $(x-1)^2 \ge 0$ και ο παρονομαστής $x^2+1 > 0$, προκύπτει ξεκάθαρα ότι:
$\boldsymbol{f'(x) \leq 0}$ για κάθε $x \in \mathbb{R}$, ενώ $f'(x) = 0$ ικανοποιείται μόνο σε μεμονωμένο σημείο ($x=1$).
Ως αποτελέσμα της κυρίαρχης αρνητικότητας και της εξασφαλισμένης συνέχειας στο $x=1$, η συνάρτηση είναι \textbf{γνησίως φθίνουσα} σε ολόκληρο το $\mathbb{R}$. 
Επειδή η $f$ μειώνεται συνεχώς δίχως να αντιστρέψει τη φορά της, \textbf{δεν συμπεριφέρει απολύτως κανένα τοπικό ούτε ολικό ακρότατο}.

\item Ζητείται $f(x) > 0$ για $x \in [0, 1]$.
Όπως αποδείξαμε, η $f$ είναι γνησίως φθίνουσα, πράγμα που σημαίνει ότι η μικρότερη τιμή της (\textbf{ελάχιστη τιμή προκειμένου το εύρος να περιορίζεται}) σε αυτό το κλειστό διάστημα λαμβάνεται αναπόφευκτα στο δεξί του άκρο $x=1$.
Υπολογίζουμε την τιμή ελαχίστου του διαστήματος:
\[ f(1) = \ln(1^2+1) - 1 + 1 = \ln 2 \]
Δεδομένου ότι το $\ln 2 \approx 0{,}69 > 0$:
\[ \text{Για κάθε } x \in [0,1], \text{ ισχύει } f(x) \geq f(1) \Rightarrow f(x) \geq \ln 2 > 0 \]
Συμπεραίνουμε πως $f(x) > 0$ για όλα τα $x \in [0,1]$.

\item Η δοθείσα εξίσωση τροποποιείται:
\[ \ln(x^2+1) = x - 1 \Leftrightarrow \ln(x^2+1) - x + 1 = 0 \Leftrightarrow f(x) = 0 \]
Χρειαζόμαστε να εξακριβώσουμε τις ρίζες της $f$ στο υποδιάστημα $\Delta = (1, +\infty)$.
- Στο αριστερό άκρο του περιορισμένου εύρους υπολογίσαμε ρητά $f(1) = \ln 2 > 0$.
- Εξετάζουμε το όριο στο $+\infty$:
\[ \lim_{x\to+\infty} f(x) = \lim_{x\to+\infty} \left( \ln(x^2+1) - x + 1 \right) = \lim_{x\to+\infty} \left[ x \left( \frac{\ln(x^2+1)}{x} - 1 + \frac{1}{x} \right) \right] \]
Συγκεκριμένα το συνοδευτικό όριο (με \eng{L'Hospital}):
\[ \lim_{x\to+\infty} \frac{\ln(x^2+1)}{x} \xlongequal[\text{\eng{DLH}}]{\frac{+\infty}{+\infty}} \lim_{x\to+\infty} \frac{\frac{2x}{x^2+1}}{1} = \lim_{x\to+\infty} \frac{2x}{x^2+1} \xlongequal[\text{\eng{DLH}}]{\dots} \lim_{x\to+\infty} \frac{2}{2x} = 0 \]
Άρα συνολικά: $(+\infty) \cdot (0 - 1 + 0) = -\infty$.

Εφόσον η $f$ είναι συνεχής και \textbf{γνησίως φθίνουσα} στο $(1, +\infty)$ και το σύνολο τιμών της είναι το $f(\Delta) = (-\infty, \ln 2)$, η τιμή $0 \in f(\Delta)$.
Δύναται άμεσα να λεχθεί ότι υπάρχει \textbf{ακριβώς και αδιαμφισβήτητα μία ρίζα} της εξίσωσης στο εύρος αυτό.

\item Διαθέτουμε κινητό σημείο $M(x(t), y(t))$ όπου $y(t) = f(x(t))$.
Δίνεται σταθερά $x'(t) = 1$. Ο ρυθμός μεταβολής της τεταγμένης επιλύεται μέσω του κανόνα αλυσίδας:
\[ y'(t) = f'(x(t)) \cdot x'(t) = f'(x(t)) \cdot 1 = f'(x(t)) \]
Τη χρονική στιγμή $t_0$ όπου το σημείο διέρχεται από την τετμημένη $x(t_0) = 1$, η τιμή αναδιπλώνεται σε:
\[ y'(t_0) = f'(1) \]
Όμως γνωρίζουμε από το (1) ότι στο $x=1$ ο αριθμητής της παραγώγου $(x-1)^2$ μηδενίζεται. Άρα $f'(1) = 0$.
Η αξιολόγηση αποδίδει άμεσα $\boldsymbol{y'(t_0) = 0 \ \text{μονάδες/χρόνο}}$.

\end{erwthma}
\end{thema}

%=========== ΛΥΣΗ - 12ο ΘΕΜΑ Δ ===========
\begin{thema}{Δ}
Δίνεται $f(x)=e^x-x-1$, $x\in\mathbb{R}$.
\begin{erwthma}

\item Παραγωγίζουμε: $f'(x) = e^x - 1$.
Επιλύουμε το πρόσημο:
- $f'(x) = 0 \Leftrightarrow e^x = 1 \Leftrightarrow x = 0$.
- $f'(x) > 0 \Leftrightarrow e^x > 1 \Leftrightarrow x > 0$.
Η συνάρτηση είναι γνησίως φθίνουσα στο $(-\infty, 0]$ και γνησίως αύξουσα στο $[0, +\infty)$.
Παρουσιάζει \textbf{ολικό ελάχιστο} στην κορυφή $x=0$, την τιμή $f(0) = e^0 - 0 - 1 = 0$.

Ως άμεση συνέπεια της απόλυτης ελαχιστοποίησης, για οποιαδήποτε τιμή $x \in \mathbb{R}$:
\[ f(x) \geq f(0) \Rightarrow e^x - x - 1 \ge 0 \Rightarrow \boldsymbol{e^x \ge x + 1} \]

\item Δίνεται $g(x) = \frac{e^x-1}{x}$ για $x\neq 0$, με $g(0)=1$.
Ελέγχουμε την ύπαρξη του ορίου στο $x=0$:
\[ \lim_{x\to 0} g(x) = \lim_{x\to 0} \frac{e^x-1}{x} \xlongequal[\text{\eng{DLH}}]{\frac{0}{0}} \lim_{x\to 0} \frac{(e^x-1)'}{(x)'} = \lim_{x\to 0} \frac{e^x}{1} = e^0 = 1 \]
Ως γνωστόν, εάν $\lim_{x\to 0} g(x) = g(0)$, η συνάρτηση είναι αυστηρώς και καθολικώς **συνεχής** στο σημείο σύνδεσης αυτό. Αφού $1 = 1$, το πόρισμα επαληθεύεται.

\item Εξετάζουμε την εφαρμογή του Θεωρήματος Μέσης Τιμής (\eng{M.V.T.}) στη συνάρτηση $f$ για το δομημένο διάστημα κόντρας $[0,1]$:
- Η $f$ είναι προφανώς συνεχής στο $[0,1]$.
- Η $f$ είναι κανονικά παραγωγίσιμη στο εσωτερικό $(0,1)$. 
Άρα εξασφαλίζεται η ύπαρξη τουλάχιστον ενός σημείου $\xi \in (0,1)$ κατά το οποίο:
\[ f'(\xi) = \frac{f(1) - f(0)}{1 - 0} \]

Υπολογίζουμε τις παραμέτρους της κλασματικής εξίσωσης:
$f(1) = e^1 - 1 - 1 = e - 2$
$f(0) = 0$ (από το Πρώτο Ερώτημα)
$f'(\xi) = e^\xi - 1$
Αντικαθιστώντας:
\[ e^\xi - 1 = \frac{e-2 - 0}{1} \Rightarrow e^\xi - 1 = e - 2 \Rightarrow e^\xi = e - 2 + 1 \Rightarrow \boldsymbol{e^\xi = e - 1} \]
Το ζητούμενο αντικρίστηκε πλήρως.

\item Θέλουμε να αποδείξουμε την ανίσωση ολοκληρωμάτων.
Γνωρίζουμε από το Ερώτημα (1) ότι $e^x \geq x + 1$ για κάθε $x \in \mathbb{R}$. 
Επιπροσθέτως, η ισότητα $e^x = x+1$ ισχύει **μόνο** στη μοναδική ρίζα $x=0$.
Επειδή οι συναρτήσεις δεν είναι ταυτόσημες στο διάστημα ολοκλήρωσης $[0,1]$, διατηρείται \textbf{αυστηρή η ανισότητα} κατά την ολοκλήρωση:
\[ \int_0^1 e^x\,\d x > \int_0^1 (x+1)\,\d x \]
Αποσυνθέτουμε και υπολογίζουμε το δεξί ολοκλήρωμα:
\[ \int_0^1 (x+1)\,\d x = \int_0^1 x\,\d x + \int_0^1 1\,\d x = \int_0^1 x\,\d x + 1 \]
Ωστόσο: $\int_0^1 x\,\d x = \left[ \frac{x^2}{2} \right]_0^1 = \frac{1}{2}$.
Οπότε:
\[ \int_0^1 (x+1)\,\d x = \frac{1}{2} + 1 = \frac{3}{2} \]
Εν τέλει, η ανασύνθεση των ταυτοτήτων αποδίδει το επιθυμητό ανισοτικό δέσιμο:
\[ \int_0^1 e^x\,\d x > 1 + \int_0^1 x\,\d x = \boldsymbol{\frac{3}{2}} \]

\end{erwthma}
\end{thema}
%=========== ΤΕΛΟΣ ΛΥΣΕΩΝ - ΔΙΑΓΩΝΙΣΜΑ ΠΡΟΣΟΜΟΙΩΣΗΣ 12 ===========
